\section{Wave-Aware Kernel Chunking}
\label{sec:optim-chunk-size}

The first \arcto{} optimization concerns the granularity at which the
batched compression and decompression kernels observe the input. The
nvCOMP-style API exposes this through the \texttt{chunk\_size}
parameter, whose inherited default is 64\,KiB. That value is a
reasonable conservative choice for pageable end-to-end execution
because it reduces metadata and bookkeeping overhead, but it is not a
kernel optimum on AMD because it does not produce enough independent
work to keep the GPU's resident wave slots populated. The relevant
architectural quantity is the resident wave-slot count
$S_{\text{wave}} = N_{\text{CU}} \cdot W_{\text{max}}$, where
$N_{\text{CU}}$ is the number of compute units and $W_{\text{max}}$
is the maximum number of resident waves per CU. On MI300X this gives
$S_{\text{wave}} = 304 \cdot 32 = 9728$ resident waves; on RX 7900 XT
the corresponding value is $84 \cdot 32 = 2688$.

The intuition behind the chunk-size choice is that the kernel sees
exactly one chunk per resident wave slot when
\[
N_{\text{chunks}} = \frac{S}{C_{\text{ker}}} = S_{\text{wave}},
\]
which gives a closed-form first-order estimate for the chunk size
that just saturates the device on a reference workload of size $S$:
\[
C_{\text{ker}}^{*} \approx \frac{S}{S_{\text{wave}}}.
\]
For a 100\,MiB reference workload, the formula predicts
$\approx 10.5\,\text{KiB}$ on MI300X, which rounds down to the nearest
power-of-two-aligned value of 8\,KiB used by the runtime. On RX 7900
XT the formula predicts $\approx 32\,\text{KiB}$. This is a
first-order model: it identifies the regime where the GPU first
becomes work-bound, not the global throughput optimum, which can lie
at smaller chunk sizes because oversubscribing the wave slots also
helps to hide memory and pipeline latency.

The empirical sweep confirms the prediction on MI300X and refines it
on the smaller GPU. Table~\ref{tab:chunk-validation} shows the
predicted $C_{\text{ker}}^{*}$ against the empirical optimum
identified by a sweep of LZ4 kernel-only compression throughput over
chunk sizes from 1\,KiB to 16\,MiB on the 100\,MiB TTI workload. The
MI300X sweep agrees with the formula at 8\,KiB and yields a
$2.87\times$ kernel-throughput speedup over the inherited 64\,KiB
default. On RX 7900 XT, the empirical optimum is one octave below the
prediction (16\,KiB rather than 32\,KiB), still well within
oversubscription range, and still well above the inherited default.
Cascaded is essentially flat across the swept range, reflecting its
more regular pipeline of delta and run-length encoding; the
formula is informative for the dictionary-style (LZ4) and
warp-cooperative (Snappy) codecs but not for Cascaded.

\begin{table}[t]
\centering
\footnotesize
\caption{Empirical validation of the wave-saturation formula
$c^{*} \approx S / W$ on the medium TTI workload ($S = 100$~MB).
The empirical optimum is taken from a sweep over chunk sizes from
$1$~KB to $16$~MB on the LZ4 compress kernel. The throughput at
default and at the empirical optimum is the kernel-side compress
throughput reported by the \arcto{} benchmark harness; the
speedup is the ratio between them.}
\label{tab:chunk-validation}
\begin{tabular}{lccccc}
\toprule
\textbf{Architecture} & \textbf{Wave slots} $W$ &
\textbf{Predicted} $c^{*}$ & \textbf{Empirical} $c^{*}$ &
\textbf{Throughput at default $64$~KB} & \textbf{Throughput at $c^{*}$} \\
\midrule
MI300X (gfx942)       & $\sim 9{,}728$ & $8$~KB  & $8$~KB  & $4.89$~GB/s     & $14.02$~GB/s ($2.87\times$) \\
RX~7900~XT (gfx1100)  & $\sim 2{,}688$ & $32$~KB & $16$~KB & $\sim 7.0$~GB/s & $\sim 12.0$~GB/s ($1.72\times$) \\
\bottomrule
\end{tabular}
\end{table}

The selected $C_{\text{ker}}^{*}$ per architecture is treated as a
fixed input to the rest of the pipeline rather than as a runtime
search. \arcto{} hardcodes the empirically validated optimum per arch
(8\,KiB on MI300X, 16\,KiB on RX 7900 XT and on the older CDNA
parts), forwards any non-zero \texttt{chunk\_size} the application
passes unchanged, and selects the architecture-aware default when
the application passes the sentinel value $0$. The end-to-end pageable
totals on a 686\,MiB MI300X input still favor 64\,KiB by a small
margin (327.96\,ms for LZ4 vs.\ a slower run at 8\,KiB) because
pageable execution is dominated by hundreds of milliseconds of H2D
and D2H rather than by the kernel; once pinned staging brings these
transfers down to roughly 14\,ms each, the kernel becomes visible
again and the smaller chunk recovers its advantage.

The two quantities, $C_{\text{ker}}$ for kernel work granularity and
$W_{\text{agg}}$ for host-device transfer granularity, are deliberately
kept decoupled: $W_{\text{agg}}$ may be tens of MiB even when
$C_{\text{ker}}$ is only a few KiB. The connection between them is
that the wave-saturation product $S_{\text{wave}} \cdot
C_{\text{ker}}^{*}$ is one of the floors that the adaptive
aggregation runtime uses to size its pinned window, as developed in
Section~\ref{sec:optim-adaptive-tiled}.
